\section{Incoming-neighborhood pattern certificates}
\label{sec:paper-incoming-certificates}

We first give a certificate that depends on a small induced subgraph but
controls second neighborhoods in the original graph. Throughout this section,
$D$ is a finite nonempty oriented graph. Put
\[
 g_D(v)=|N_D^+(v)|-|N_D^{++}(v)|,
\]
where $N_D^{++}(v)$ consists of the vertices at directed distance exactly
two from $v$. A vertex with $g_D(v)\leq0$ is called a Seymour vertex.
No extremal or minimal-counterexample assumption is made in this section.

\subsection{Patterns and the transfer theorem}

Let $H$ be a finite nonempty oriented graph. Define its family of
\emph{permitted incoming patterns} by
\begin{equation}
 \mathcal P(H)=\{P\subseteq V(H):
           P\subseteq N_H^-(j)\text{ for some }j\in V(H)\}.
 \label{eq:paper-pattern-family}
\end{equation}
In particular, the empty set belongs to $\mathcal P(H)$. For an arbitrary
$P\subseteq V(H)$, write
\[
 R_H(P)=\{i\in V(H)\setminus P:
                      N_H^+(i)\cap P\ne\varnothing\}.
\]
Thus $R_H(P)$ consists of the vertices outside $P$ that send an arc into
$P$. If $\lambda$ is a nonnegative vector indexed by $V(H)$, use the
notation $\lambda(S)=\sum_{i\in S}\lambda_i$.

\begin{definition}[Pattern certificate]
\label{def:paper-pattern-certificate}
A pattern certificate for $H$ is a nonzero vector $\lambda\geq0$ such that
\begin{equation}
 \lambda(R_H(P))\geq\lambda(P)
             \qquad\text{for every }P\in\mathcal P(H).
 \label{eq:paper-pattern-inequalities}
\end{equation}
The inequalities are homogeneous, so a certificate may always be
normalized by $\lambda(V(H))=1$.
\end{definition}

For a nonempty set $C\subseteq V(D)$, let $H=D[C]$ and associate to every
original vertex $v$ the pattern
\[
 \Pi_C(v)=N_D^-(v)\cap C.
\]
We say that $C$ is \emph{incoming-covered} if
\begin{equation}
 \Pi_C(v)\in\mathcal P(H)\qquad(v\in V(D)).
 \label{eq:paper-incoming-covered}
\end{equation}
The witness for this condition may depend on $v$; the condition does not
require any particular orientation between $v$ and that witness.

For clarity, let $\mathbf N_D$ denote the matrix with rows and columns
indexed by $V(D)$ and entries
\[
 (\mathbf N_D)_{uv}=
 \begin{cases}
 -1,&v\in N_D^+(u),\\
  1,&v\in N_D^{++}(u),\\
  0,&\text{otherwise}.
 \end{cases}
\]
The row sum at $u$ is $-g_D(u)$.

\begin{theorem}[Transfer from an incoming-covered core]
\label{thm:pattern-transfer}
Let $C\subseteq V(D)$ be nonempty and incoming-covered, and put $H=D[C]$.
If $\lambda$ is a pattern certificate for $H$, then its extension
$\widehat\lambda$ by zero to $V(D)$ satisfies
\[
 \mathbf N_D^{\mathsf T}\widehat\lambda\geq0.
\]
Consequently some $c\in C$ with $\lambda_c>0$ is a Seymour vertex of $D$.
\end{theorem}
\begin{proof}
Fix an original target $v$ and put $P=\Pi_C(v)$. The first predecessors of
$v$ in $C$ are precisely the vertices of $P$. If $i\in R_H(P)$, choose
$j\in N_H^+(i)\cap P$. Both arcs of the path $i\to j\to v$ are
original arcs of $D$, and $i\notin P$ says that $i\to v$ is not an arc.
Also $i\ne v$: otherwise $v\to j$ and $j\to v$ would form a digon.
It follows that $v\in N_D^{++}(i)$.

The total $\lambda$-mass of the core vertices that have $v$ as an exact
second neighbor is therefore at least $\lambda(R_H(P))$. Since $P$ is
permitted, this is at least $\lambda(P)$, the total mass of its core
first predecessors. This proves the asserted matrix inequality at the
column indexed by $v$.

Summing over the columns gives
\[
 0\leq\mathbf1^{\mathsf T}\mathbf N_D^{\mathsf T}\widehat\lambda
       =-\sum_{c\in C}\lambda_c g_D(c).
\]
Since $\lambda$ is nonzero and nonnegative, not every vertex of its
support can have positive gap.
\end{proof}

The certificate weights the possible sources in $C$. The neighborhoods
and their cardinalities in the conclusion remain those of the unweighted
original graph. In particular, vertices outside $C$ are retained as
targets in every column of the transfer inequality.

\subsection{Tournament cores}

Tournament cores admit certificates of arbitrary order. We use the
losing-density formulation underlying Fisher's proof for tournaments
\cite{Fisher1996}. The short argument below also specifies the sign
convention for the skew matrix.

\begin{lemma}[A skew-game density]
\label{lem:paper-skew-density}
Let $B$ be a real skew-symmetric matrix. There is a nonnegative vector
$\lambda$ with $\mathbf1^{\mathsf T}\lambda=1$ and $B\lambda\geq0$.
\end{lemma}
\begin{proof}
Let $\Delta$ be the probability simplex of the appropriate dimension.
Finite minimax gives
\[
 \max_{\lambda\in\Delta}\min_{p\in\Delta}p^{\mathsf T}B\lambda
 =\min_{p\in\Delta}\max_{\lambda\in\Delta}p^{\mathsf T}B\lambda.
\]
The left-hand side is at most zero, by choosing $p=\lambda$ in the
inner minimum. The right-hand side is at least zero, by choosing
$\lambda=p$ in the inner maximum. Both values are therefore zero.
A maximizing $\lambda$ satisfies $p^{\mathsf T}B\lambda\geq0$ for
every $p\in\Delta$, and in particular for every coordinate vector.
\end{proof}

If $A_H$ is the adjacency matrix of a tournament $H$, put
$B_H=A_H-A_H^{\mathsf T}$. A vector supplied by
Lemma~\ref{lem:paper-skew-density} is called a losing density of $H$.
At each vertex, its outgoing density is at least its incoming density.

\begin{lemma}[Expansion of covered patterns in tournaments]
\label{lem:paper-tournament-pattern-expansion}
Every losing density of a nonempty tournament $H$ is a pattern certificate
for $H$.
\end{lemma}
\begin{proof}
Fix $P\subseteq N_H^-(j)$ and set
\[
 R=R_H(P),\qquad F=V(H)\setminus(P\cup R).
\]
The witness $j$ belongs to $F$, so $F$ is nonempty. Tournament
completeness implies that every arc between $P$ and $F$ is directed from
$P$ to $F$: a vertex of $F$ has no arc into $P$ by its definition.

First suppose that $\lambda(F)>0$. Sum the nonnegative quantities
$\lambda_i(B_H\lambda)_i$ over $i\in F$. The terms internal to $F$
cancel by skew-symmetry. The contribution from pairs with the other
endpoint in $P$ is $-\lambda(F)\lambda(P)$. The contribution from pairs
with the other endpoint in $R$ is at most $\lambda(F)\lambda(R)$.
Consequently
\[
 0\leq\sum_{i\in F}\lambda_i(B_H\lambda)_i
 \leq\lambda(F)\bigl(\lambda(R)-\lambda(P)\bigr),
\]
and division by $\lambda(F)$ proves the required inequality.

If $\lambda(F)=0$, apply $(B_H\lambda)_j\geq0$ at any $j\in F$.
The incoming mass at $j$ is at least $\lambda(P)$, whereas its outgoing
mass is at most $\lambda(R)$. This again gives
$\lambda(R)\geq\lambda(P)$. The argument also covers $P=\varnothing$.
\end{proof}

\begin{corollary}[Incoming-covered tournament cores]
\label{cor:paper-tournament-core}
If a nonempty incoming-covered set $C$ induces a tournament in $D$, then
$C$ contains a Seymour vertex of $D$. More precisely, the zero extension
of every losing density of $D[C]$ satisfies the matrix inequality in
Theorem~\ref{thm:pattern-transfer}.
\end{corollary}
\begin{proof}
Combine Lemmas~\ref{lem:paper-skew-density}
and~\ref{lem:paper-tournament-pattern-expansion} with
Theorem~\ref{thm:pattern-transfer}.
\end{proof}

\subsection{Maximal incoming classes and exact small-core certificates}

Define an equivalence relation on $V(D)$ by
\[
 u\sim_D v\quad\Longleftrightarrow\quad N_D^-(u)=N_D^-(v).
\]
The equivalence classes are partially ordered by containment of their
common incoming sets. We call them \emph{incoming classes}, and
maximality will always refer to this inclusion order. Distinct members
of one class cannot be adjacent, since an internal arc would put its tail
in its own incoming set by equality of the two incoming neighborhoods.
An incoming class need not have a common outgoing neighborhood.

\begin{lemma}[Representative cores are incoming-covered]
\label{lem:paper-maximal-representatives-covered}
Choose one representative from each maximal incoming class of $D$, and
let $C$ be the set of chosen vertices. Then $C$ is incoming-covered.
\end{lemma}
\begin{proof}
For every $v\in V(D)$, finiteness of the inclusion poset gives a maximal
class with representative $c$ such that
$N_D^-(v)\subseteq N_D^-(c)$. Intersecting with $C$ gives
\[
 \Pi_C(v)\subseteq N_D^-(c)\cap C=N_{D[C]}^-(c).
\]
This is exactly the required condition.
\end{proof}

\begin{lemma}[Exact certificates for cores of order at most five]
\label{lem:paper-exact-five-certificates}
Every nonempty oriented graph of order at most five has a pattern
certificate. A certificate can be chosen with integer entries in
$\{0,1,3\}$.
\end{lemma}
\begin{proof}[Computer-assisted proof]
The certificate data consist of one nonzero nonnegative integer vector
for every labelled oriented graph of orders one through five. The
encoding lists the unordered pairs
\[
 (0,1),(0,2),\ldots,(0,n-1),(1,2),\ldots,(n-2,n-1)
\]
in lexicographic order. The corresponding ternary digit is zero for a
missing pair, one for the arc from the smaller endpoint to the larger
endpoint, and two for the reverse arc. The first pair corresponds to
the least significant digit. Thus the codes
$0,\ldots,3^{\binom n2}-1$ cover every labelled oriented graph exactly
once.

For every code, the verifier reconstructs its adjacency matrix and checks
that the supplied vector has length $n$, nonnegative integer entries, and
positive sum. It then enumerates every subset $P\subseteq\{0,\ldots,n-1\}$.
A subset is retained precisely when some $j$ satisfies $i\to j$ for
every $i\in P$. For each retained subset, the verifier forms
$R_H(P)$ directly from the reconstructed adjacency matrix and checks
\eqref{eq:paper-pattern-inequalities} with integer arithmetic.

The numbers of graph codes for $n=1,2,3,4,5$ are, respectively,
\[
 1,\quad3,\quad27,\quad729,\quad59049.
\]
All $59809$ certificates pass all permitted-pattern inequalities. There
are $522213$ such checks, of which $462404$ use a nonempty pattern.
The largest integer entry is $3$, and the largest sum of entries is $9$.
Among the order-five certificates, $59009$ are uniform on their supports
and the remaining $40$ are nonuniform. This exhausts the finite assertion.

The complete tables and verification sources accompany the manuscript.
In particular, \texttt{verify\_core\_pattern\_certificates\_independent.js}
reconstructs Boolean adjacency matrices and checks the inequalities
without calling an optimization solver or using a generator's graph
utilities. The integer tables, rather than a floating-point optimization
result, are the certificates used in this proof.
\end{proof}

\begin{theorem}[An entire Seymour maximal class]
\label{thm:core-five}
Every finite nonempty oriented graph with at most five maximal incoming
classes has a maximal incoming class all of whose vertices are Seymour
vertices. There is no restriction on the total number of vertices or on
the sizes of the classes.
\end{theorem}
\begin{proof}
Suppose that no maximal incoming class consists entirely of Seymour
vertices. Choose from each maximal class a vertex with positive gap,
and let $C$ be the set of chosen representatives. Then $1\leq|C|\leq5$,
and Lemma~\ref{lem:paper-maximal-representatives-covered} shows that
$C$ is incoming-covered. Lemma~\ref{lem:paper-exact-five-certificates}
supplies a pattern certificate for $D[C]$.
Theorem~\ref{thm:pattern-transfer} consequently gives a Seymour vertex
in $C$, contrary to the selection of all its members.
\end{proof}

\begin{corollary}
\label{cor:paper-counterexample-six-incoming-classes}
Every counterexample to Seymour's second neighborhood conjecture has at
least six maximal incoming classes.
\end{corollary}
\begin{proof}
An all-deficient graph has no Seymour vertex, so
Theorem~\ref{thm:core-five} applies contrapositively.
\end{proof}

The role of the computation is confined to the finitely many cores in
Lemma~\ref{lem:paper-exact-five-certificates}. The passage to arbitrary
orders and arbitrary class sizes in Theorem~\ref{thm:core-five} uses the
original-graph transfer theorem. It is not an enumeration of all graphs
to which that theorem applies.

\subsection{Strict duality and realization of an infeasible core}

The pattern system also gives an equivalent formulation of the full
conjecture. Its reverse implication requires a realization argument:
infeasibility of a system on a core must yield an actual oriented graph
with positive gap at every vertex.

\begin{theorem}[Pattern feasibility and the second neighborhood conjecture]
\label{thm:pattern-duality}
The following statements are equivalent:
\begin{enumerate}
 \item Every finite nonempty oriented graph has a Seymour vertex.
 \item Every finite nonempty oriented graph $H$ has a pattern certificate.
\end{enumerate}
More explicitly, if the normalized system
\[
 \lambda\geq0,\qquad \sum_i\lambda_i=1,\qquad
 \lambda(R_H(P))\geq\lambda(P)\quad(P\in\mathcal P(H))
\]
is infeasible for some $H$, a finite all-deficient oriented graph can be
constructed from $H$ and a rational strict dual solution.
\end{theorem}
\begin{proof}
First assume statement~(2). For an arbitrary $D$, take the core to be
$D$ itself. Its pattern at a target $v$ is the full set $N_D^-(v)$,
which is permitted. Theorem~\ref{thm:pattern-transfer} yields a Seymour
vertex, proving statement~(1).

For the converse, suppose that the normalized system for a nonempty
core $H$ is infeasible. For each permitted pattern, define the column
vector
\[
 a_P(i)=\mathbf1_{i\in R_H(P)}-\mathbf1_{i\in P}.
\]
On the probability simplex $\Delta$ indexed by $V(H)$, the continuous
function $\min_{P\in\mathcal P(H)}\lambda^{\mathsf T}a_P$ is strictly
negative everywhere. Compactness therefore gives
\begin{equation}
 \max_{\lambda\in\Delta}
       \min_{P\in\mathcal P(H)}\lambda^{\mathsf T}a_P<0.
 \label{eq:paper-strict-dual-value}
\end{equation}
By finite minimax, there are nonnegative real pattern weights whose
weighted sum of the $a_P$ has every coordinate strictly negative.
They may be chosen rational: keep their zero coordinates equal to zero
and approximate the positive coordinates sufficiently closely by
positive rationals, using strictness of all the finitely many
inequalities. Clearing denominators gives nonnegative integers $w_P$
such that
\begin{equation}
 \gamma_i=\sum_{P\in\mathcal P(H)}w_P
       \bigl(\mathbf1_{i\in P}-\mathbf1_{i\in R_H(P)}\bigr)>0
       \qquad(i\in V(H)).
 \label{eq:paper-positive-dual-gaps}
\end{equation}

For each permitted pattern $P$, choose a witness $j(P)$ with
$P\subseteq N_H^-(j(P))$. Fix a positive integer $K$. The new graph
$D_K$ has two kinds of vertices. For each $j\in V(H)$, it has a
distinguished vertex $c_j$ with label $j$ and pattern $N_H^-(j)$.
For each $P\in\mathcal P(H)$, it also has $K w_P$ vertices with
pattern $P$ and label $j(P)$. Distinct vertices are retained when
patterns or labels coincide. If $\ell(v)$ and $P(v)$ denote the label
and pattern of a vertex, define
\begin{equation}
 v\to u\quad\Longleftrightarrow\quad \ell(v)\in P(u).
 \label{eq:paper-pattern-realization-arcs}
\end{equation}

Every pattern satisfies $P(v)\subseteq N_H^-(\ell(v))$, so
$\ell(v)\notin P(v)$; hence \eqref{eq:paper-pattern-realization-arcs}
creates no loop. If both $v\to u$ and $u\to v$ held, the same
containments would give both label arcs $\ell(v)\to\ell(u)$ and
$\ell(u)\to\ell(v)$ in $H$. For unequal labels this is a digon,
and for equal labels it is a loop. Thus $D_K$ is oriented. Its
distinguished vertices induce exactly $H$.

Consider any source $v$ with label $i$. Its first targets are exactly
the vertices $u$ with $i\in P(u)$. We claim that its exact second
targets are exactly those with $i\in R_H(P(u))$. Indeed, if
$v\to x\to u$ and $i\notin P(u)$, put $j=\ell(x)$. The first
arc and $P(x)\subseteq N_H^-(j)$ imply $i\to j$ in $H$, while
the second arc says $j\in P(u)$. Thus $i\in R_H(P(u))$.
Conversely, if $i\in R_H(P(u))$, choose
$j\in N_H^+(i)\cap P(u)$. The distinguished vertex $c_j$ supplies
the path $v\to c_j\to u$, and $i\notin P(u)$ excludes a first
arc. The target cannot be $v$: since
$P(v)\subseteq N_H^-(i)$, such a $j$ would give a digon $i\leftrightarrow j$
in $H$.

The extra pattern vertices therefore contribute $K\gamma_i$ to the
gap of every source labelled $i$. The distinguished vertices contribute
$g_H(i)$, because they supply exactly the first and exact second
targets of $i$ in $H$. We have proved the identity
\begin{equation}
 g_{D_K}(v)=K\gamma_i+g_H(i)
                    \qquad\text{whenever }\ell(v)=i.
 \label{eq:paper-realization-gap}
\end{equation}
Choose $K$ so large that the right-hand side is positive for every
$i\in V(H)$. This is possible by
\eqref{eq:paper-positive-dual-gaps}. Then $D_K$ is a finite
all-deficient oriented graph. Consequently statement~(1) rules out
infeasibility for every $H$, proving statement~(2).
\end{proof}

\begin{corollary}[An equivalent whole-class formulation]
\label{cor:paper-whole-class-equivalence}
Seymour's second neighborhood conjecture is equivalent to the assertion
that every finite nonempty oriented graph has a maximal incoming class
consisting entirely of Seymour vertices.
\end{corollary}
\begin{proof}
The whole-class assertion plainly implies the conjecture. Conversely,
assume the conjecture, and hence universal pattern feasibility by
Theorem~\ref{thm:pattern-duality}. If each maximal incoming class of
$D$ contained a positive-gap vertex, choose one such representative
from each class. Lemma~\ref{lem:paper-maximal-representatives-covered}
and Theorem~\ref{thm:pattern-transfer} would give a Seymour vertex
among these representatives, a contradiction.
\end{proof}

\subsection{Lifting through iterated representative cores}

The representative operation can be repeated, and it preserves more
than the existence of one Seymour vertex: a certificate at a later
stage lifts to a certificate at every preceding stage.

\begin{lemma}[Lifting a pattern certificate]
\label{lem:core-lifting}
Let $C$ contain one representative from each maximal incoming class of
an oriented graph $H$, and put $K=H[C]$. The extension by zero of any
pattern certificate for $K$ is a pattern certificate for $H$.
\end{lemma}
\begin{proof}
Fix $P\subseteq N_H^-(v)$. Choose a representative $c\in C$ such
that $N_H^-(v)\subseteq N_H^-(c)$. Then
\[
 P\cap C\subseteq N_K^-(c),
 \qquad R_K(P\cap C)\subseteq R_H(P)\cap C.
\]
For the second inclusion, a vertex of $C$ outside $P\cap C$ is
outside $P$, and an arc into $P\cap C$ is also an arc of $H$ into
$P$. If $\lambda$ is the certificate on $C$ and
$\widehat\lambda$ is its zero extension, it follows that
\[
 \widehat\lambda(R_H(P))
 \geq\lambda(R_K(P\cap C))
 \geq\lambda(P\cap C)
 =\widehat\lambda(P).
\]
The lifted vector is nonnegative and nonzero, so it is a certificate.
\end{proof}

\begin{corollary}[Iterated-core sufficient conditions]
\label{cor:paper-iterated-core-sufficient}
Starting from $D$, repeatedly take the induced subgraph on one
representative from each maximal incoming class of the current graph.
If any resulting core has at most five vertices, or is a tournament
of any order, then the original graph $D$ has a Seymour vertex in that
core.
\end{corollary}
\begin{proof}
The small-core certificate lemma or the tournament certificate lemma
supplies a certificate at the specified stage. Apply
Lemma~\ref{lem:core-lifting} successively in reverse order. This gives
a certificate on $D$ supported in the specified core. Applying
Theorem~\ref{thm:pattern-transfer} with the whole graph as its core
gives an original Seymour vertex in that support.
\end{proof}

\begin{corollary}[Reduction to incomparable incoming neighborhoods]
\label{cor:paper-terminal-incoming-core}
Universal pattern feasibility, and therefore Seymour's second neighborhood
conjecture, is equivalent to pattern feasibility for oriented graphs
whose incoming neighborhoods are distinct and pairwise incomparable
under inclusion.
\end{corollary}
\begin{proof}
Whenever the representative operation reduces the vertex set, repeat
it. The process terminates because the order decreases. If the
operation no longer reduces the order, every incoming class is a
singleton and every class is maximal. The incoming neighborhoods are
therefore distinct and pairwise incomparable. A certificate for this
terminal graph lifts to the original graph by
Lemma~\ref{lem:core-lifting}. The reverse implication is immediate.
\end{proof}

This is an induced-subgraph reduction: it deletes vertices and does not
add arcs. Its conclusion concerns the original neighborhoods of vertices
in the final core. For a single chosen chain it does not assert that
an entire original maximal incoming class consists of Seymour vertices;
the whole-class conclusions above use representative selection under
the negation of that stronger assertion.
